	\documentclass[fleqn, 12pt]{article}
	\usepackage[latin1]{inputenc}
	\usepackage[ngerman]{babel}
	\usepackage[a4paper,text={155mm,220mm},centering,headsep=10mm,footskip=15mm]{geometry}
	\usepackage{graphicx}
	\usepackage{amssymb}
	\usepackage{graphicx}
	\usepackage{multirow}
	\usepackage{multicol}
	\usepackage{pst-plot}
	\usepackage{float}
	\usepackage{pdfpages}
	\title{}
	\date{}
	\author{}
	\usepackage{array}
	\usepackage{amsmath}
	\usepackage{fancyhdr}
	\usepackage{enumitem} 
	
	% Listings
	\usepackage{listings}
	\renewcommand{\lstlistlistingname}{Liste der Quellcode Ausschnitte}
	\renewcommand{\lstlistingname}{Quellcode Ausschnitt}
		\definecolor{lsgreen}{rgb}{0,.5,0}
		\definecolor{lsred}{rgb}{.7,0,0}
		\definecolor{lsorange}{rgb}{.9,.5,0}
		\definecolor{lsgray}{rgb}{.5,.5,.5}
		\lstset{
			frame=tb,
			aboveskip=10mm,
			belowskip=10mm,
			showstringspaces=false,
			columns=flexible,
			captionpos=b,
			basicstyle={\normalsize\ttfamily},
			numbers=left,
			numberstyle=\tiny\color{lsgray},
			keywordstyle=\color{blue},
			commentstyle=\color{lsorange},
			stringstyle=\color{lsgreen},
			breaklines=true,
			breakatwhitespace=true,
			rulecolor=\color{lsgray},
			xleftmargin=7mm,
			tabsize=3
		}
	
	% Needed for proper display of links in bibliography
	\usepackage{url}	
	
	% Hyperref
	\usepackage[
	pdfauthor={Laura Anger, Timo \dots, Lukas Kolhagen},
	pdftitle={BVA Bewegungsanalyse},
	pdftoolbar=true,	
	colorlinks=true,
	linkcolor=blue,
	citecolor=blue,
	urlcolor=blue,
	linktocpage=true
	]{hyperref}
	\usepackage{bookmark}
	\bookmarksetup{
	numbered
	}
	\urlstyle{same}
	
	\pagestyle{fancy}
	\fancyhf{}
	\fancyhead[LO]{BVA}
	\fancyhead[CO]{Bewegungsanalyse in einer Videosequenz}
	\fancyhead[RO]{\thepage}
	\renewcommand{\headrulewidth}{0.5pt}
	\newcommand{\Absatzbox}[1]{\parbox[0pt][2em][c]{0cm}{}}
	
	% Itemize symbol
	\renewcommand{\labelitemi}{$\triangleright$}
	
	% Colors
	\definecolor{yellow}{rgb}{.95,.85,0}
	\definecolor{green}{rgb}{0,.8,0}
	\definecolor{blue}{rgb}{0,0,.8}
	\definecolor{red}{rgb}{.8,0,0}
	\definecolor{grey}{rgb}{.4,.4,.4}
	\definecolor{orange}{rgb}{.9,.5,0}
	
	% Macro for typesetting C++
	\newcommand{\CC}{C\nolinebreak\hspace{-.05em}\raisebox{.4ex}{\tiny\bf +}\nolinebreak\hspace{-.10em}\raisebox{.4ex}{\tiny\bf +}}
\def\CC{{C\nolinebreak[4]\hspace{-.05em}\raisebox{.4ex}{\tiny\bf ++}}}

\begin{document}
\thispagestyle{empty}
	%\sffamily
			\begin{center}
			\includegraphics[width=.35\textwidth]{logo_TH}\\[20ex]
			{\Huge\textbf{Projektdokumentation}}\\[8ex]
			\rule{.8\textwidth}{.2pt}
			{\Large Bewegungsanalyse in einer Videosequenz\\[1ex] mit dem Ansatz des Papers
			von Aach und Kunz}\\
			\rule{.8\textwidth}{.2pt}\\[10ex]
			von\\[2ex]
			\begin{tabular}{ll}
			Laura Anger &(Matrikelnr. XXXXXXXX)\\ 
			Timo \dots &(Matrikelnr. XXXXXXXX)\\ 
			Lukas Kolhagen &(Matrikelnr. 11084355)\\
			\end{tabular}\\[10ex]
			Durchgef�hrt im\\ \textbf{Master Medientechnologie}\\ im\\ 
			\textbf{Sommersemester 2016}\\			
			\end{center}
			\vfill
			\begin{flushleft}
			{\bf Betreuer:}\\
			Prof. Dr. Dietmar Kunz\\
			Institut f�r Medien- und Phototechnik
			\end{flushleft}
	\newpage
	\tableofcontents
	\newpage
	\section{Einleitung}
	Diese Ausarbeitung ist Teil der Abschlussprojekt-Dokumentation im Modul "`Weiterf�hrende Themen der Bildverarbeitung"' im Master Medientechnologie an der Technischen Hochschule K�ln.
	
	Das Projekt besch�ftigte sich mit der Bewegungsanalyse einer Videosequenz mit dem Ansatz des Papers von Aach und Kunz\textsuperscript{\cite{aach1998bayesian}}. Es wurde bearbeitet von Laura Anger, Timo \dots und Lukas Kolhagen.
		\subsection{Ansatz im Paper von Aach und Kunz}
	
	\section{Implementierung}
		\subsection{Framework f�r die Darstellung}
		\subsection{Bewegungssch�tzung}
			\subsubsection{Bildrauschen}
			\subsubsection{�rtliche Koh�renz}
			\subsubsection{Zeitliche Koh�renz}
			\subsubsection{Zu minimierende Kostenfunktion}
	\section{Auswertung}
		\subsection{Testvideos}
		\subsection{Ergebnisse}
	\section{Zusammenfassung}
	\newpage
	\bibliographystyle{plain}
	\bibliography{ABV_Bewegungsanalyse_LaTeX}
\end{document}
		